% Generated draft: requires mathematical and source review.
\documentclass[11pt]{article}
\usepackage[margin=25mm]{geometry}
\usepackage{fontspec}
\setmainfont{Latin Modern Roman}
\usepackage{amsmath,amssymb,mathtools,unicode-math}
\setmathfont{Latin Modern Math}
\usepackage{enumitem,needspace,xurl,hyperref}
\setlength{\parindent}{0pt}
\setlength{\parskip}{5pt}
\setlength{\emergencystretch}{3em}
\newcommand{\N}{\mathbb{N}}
\newcommand{\Z}{\mathbb{Z}}
\newcommand{\Q}{\mathbb{Q}}
\newcommand{\R}{\mathbb{R}}
\newcommand{\C}{\mathbb{C}}
\newcommand{\F}{\mathbb{F}}
\newcommand{\A}{\mathbb{A}}
\newcommand{\PP}{\mathbb{P}}
\newcommand{\E}{\mathbb{E}}
\newcommand{\eps}{\varepsilon}
\newcommand{\dd}{\,\mathrm{d}}
\begin{document}
\section*{UnsolvedMath ID: 30000388}
\textbf{Source problem code:} catalogue entry
\subsection{Definitions and mathematical statement}
Let $S^d\subset\R^{d+1}$ be the unit sphere and let $\mathcal A=(H_1,\ldots,H_n)$ be an arrangement of great hemispheres. Using the closed-hemisphere convention, write
\[
 H_i=\{x\in S^d:\langle a_i,x\rangle\ge0\},\qquad a_i\ne0.
\]
Vertices are the zero-dimensional cells of the arrangement of the bounding great spheres, counted as distinct points. A vertex's level is
\[
 \operatorname{lev}_{\mathcal A}(x)
 =\#\{i:x\notin H_i\}=\#\{i:\langle a_i,x\rangle<0\}.
\]
Thus a hemisphere whose boundary contains the vertex does contain it for this convention. Put
\[
 V_{\le\ell}(\mathcal A)
 =\#\{x\text{ a vertex of }\mathcal A:
                     \operatorname{lev}_{\mathcal A}(x)\le\ell\}.
\]
The conjecture asks whether
\[
 V_{\le\ell}(\mathcal A)\le v(\ell,n,d)
 \qquad\text{for every integer }0\le\ell<(n-d)/2,
\]
where $v(\ell,n,d)$ is the corresponding count for the arrangement polar-dual to a cyclic $d$-polytope with $n$ vertices, as specified by the report. A cyclic polytope is the convex hull of $n$ points on the moment curve $(t,t^2,\ldots,t^d)$. The comparison concerns all levels up to $\ell$, not only vertices at exactly level $\ell$.

The extract does not specify general-position/simplicity assumptions, degeneracy conventions or the detailed homogeneous polar construction of the benchmark. These must be matched to the source rather than assigning an unsupported explicit formula for $v$. The closed-hemisphere convention is stated because switching to open hemispheres would change levels of boundary vertices.

The literature assessment checked on 21 August 2026 calls this a generalized upper-bound conjecture and reports positive low-dimensional and special cases, including the case $n=d+4$ in 2025 work. It says the all-dimension statement remains open. No general conclusion is inferred from that special vertex-count regime; the assessment is dated supplied evidence requiring independent verification.
\subsection{Short English statement}
Are low-level vertex counts in any great-hemisphere arrangement bounded by those of the cyclic-polytope polar model? The exact range is $0\le\ell<(n-d)/2$, and all vertices of level at most $\ell$ are counted.
\subsection{Sources}
\begin{itemize}
\item \emph{Combinatorics} (2006), Oberwolfach Reports, DOI 10.4171/owr/2006/01, pp. 5--72, \url{https://doi.org/10.4171/owr/2006/01}.
\item Literature assessment checked 21 August 2026: ``Levels in arrangements,'' \url{https://doi.org/10.4230/LIPIcs.SoCG.2025.75}; original report, \url{https://doi.org/10.4171/owr/2006/01}.
\item Local UnsolvedMath registry, permanent ID $30000388$, exact source code \texttt{catalogue entry}, created 25 August 2026 and updated 16 September 2026. The clean statement is labelled corrected and verified against the report, retaining its vertex-level definition and cyclic benchmark.
\end{itemize}
\end{document}
